\conclusions \label{sec:conclusions}
%
% Short summary
%
This study explored the potential of a hybrid strategy for volcanic-cloud modelling based on the integration of generative AI with physics-based simulations. A Variational Autoencoder (VAE) was trained on a relatively small ensemble of numerical simulations to learn a compact, low-dimensional representation of the high-dimensional problem. This methodology offers several potential advantages, including the generation of large ensembles at very low computational cost, the reduction of sampling errors and improved estimation of ensemble statistics, the generation of new samples beyond those contained in the training dataset, and the ability to perform statistical operations in a substantially lower-dimensional space, including latent-space data assimilation.

%
% Main results
%
Two numerical experiments have been presented. A first synthetic experiment assessed the consistency of the VAE-generated ensemble using a controlled case characterized by multiple qualitatively different solutions, testing the ability of the VAE to capture this variability. Even when no new physical knowledge is introduced by a VAE-augmented ensemble, it provides a good statistical approximation of the reference distribution and generates fields that are statistically consistent with the underlying physical model, particularly at large spatial scales. The results also highlight limitations in the representation of smaller-scale spatial structures.
%
However, the reconstruction error is smaller than the errors introduced by the numerical model itself. In the second  experiment, the 24 December 2018 Etna eruption was modelled and the VAE-augmented ensemble was compared with satellite retrievals. The agreement with observations was found very similar for the numerical and VAE-augmented ensembles, with the latter yielding better RMSE, RMSLE, and Pearson correlation coefficients, suggesting that the generation of new samples can lead to a better statistical representation of forecast uncertainty.
%
The Etna experiment was also used to provide a proof-of-concept demonstration of latent-space data assimilation methodologies. A two-dimensional estimation problem was solved by assimilating satellite retrievals in the VAE latent space. The analysis ensemble showed improved agreement with the observations across all evaluated metrics, while significantly reducing the ensemble spread and retaining the ability to represent the spatial variability of the volcanic cloud and the physical consistency with the simulated prior. 

%
% Final summary
%
In summary, this work represents a first proof of concept based on a relatively simple neural architecture and an elementary data assimilation methodology. The results demonstrate the potential of the formulation, while highlighting the need for further research to assess its robustness and applicability across different eruption scenarios. Future work should extend the methodology to three-dimensional fields, explore more advanced generative models and assimilation approaches, and evaluate their performance against state-of-the-art DA methods and independent observations, ultimately assessing their potential implications for operational forecasting.